\chapter{Risultati}
\label{cap:4}

Questo capitolo presenta i risultati della valutazione comparativa tra
il modello addestrato con Reflection Tuning e quello addestrato con
Knowledge Distillation response-based. La valutazione è stata condotta
su un insieme di 20 domande di ragionamento matematico, giudicate
da Gemini Flash secondo il paradigma LLM-as-a-Judge descritto nel
capitolo precedente. Per ciascuna domanda sono stati raccolti punteggi
su tre metriche (Correctness, Reasoning Quality, Clarity) e un
giudizio complessivo di vittoria.


\section{Risultati Quantitativi}
\label{sec:risultati_quantitativi}

\subsection{Pipeline sperimentale e dati raccolti}
\label{subsec:pipeline_risultati}

Prima di presentare i risultati numerici è utile richiamare la
struttura della pipeline sperimentale, illustrata nella
Figura~\ref{fig:pipeline} del capitolo precedente.
Il punto di partenza è il dataset pubblico \textbf{OpenThoughts-114k},
composto da circa 114\,000 esempi di ragionamento generati da modelli
di grandi dimensioni. Da questo corpus sono stati estratti e filtrati
\textbf{700 campioni di training} e \textbf{80 campioni di valutazione},
selezionando conversazioni con tracce di ragionamento di lunghezza e
qualità sufficienti a garantire la supervisione del modello.

A partire dagli stessi 700 campioni, la pipeline si biforca in due
rami paralleli e simmetrici, ciascuno destinato ad addestrare una
versione distinta di Llama~3.2~1B-Instruct. Nel \textbf{ramo
Reflection}, le 700 conversazioni originali --- già dotate di un
blocco \texttt{<thinking>} con il ragionamento esplicito generato
da modelli di grandi dimensioni --- vengono utilizzate direttamente
come target di fine-tuning: lo student impara a produrre, per ogni
domanda, sia la traccia di ragionamento sia la risposta finale,
replicando la struttura a due stadi del dato originale. Nel \textbf{ramo
Distillation}, le stesse 700 domande vengono sottoposte a Gemini
Flash (il modello \emph{teacher}), che restituisce per ciascuna una
risposta diretta e concisa, priva di ragionamento esplicito: questo
secondo dataset di 700 coppie domanda--risposta viene usato per
addestrare lo student secondo il paradigma del \emph{behavior cloning},
in cui l'obiettivo è imitare lo stile e la correttezza del teacher.

Entrambi i modelli vengono addestrati con \textbf{QLoRA} (rango
$r = 16$, $\alpha = 32$, quantizzazione 4-bit) per lo stesso numero
di epoche e con gli stessi iperparametri, in modo da rendere il
confronto finale attribuibile esclusivamente alla natura del dato
di training e non a differenze nella procedura di ottimizzazione.
Al termine del training, i due modelli vengono valutati su un test
set comune di \textbf{20 domande di ragionamento matematico},
somministrate a entrambi nello stesso ordine. Le risposte vengono
giudicate da Gemini Flash secondo il paradigma \textbf{LLM-as-a-Judge}
\cite{zheng2023judging}: per ogni domanda il giudice assegna un
punteggio da 1 a 5 su tre metriche (\emph{Correctness},
\emph{Reasoning Quality}, \emph{Clarity}) e un giudizio complessivo
di vittoria. La domanda Q18 ha prodotto un errore API durante la
fase di evaluation e non è inclusa nel calcolo delle medie ($n = 19$).

La Tabella~\ref{tab:medie} riassume i punteggi medi per metrica;
la Tabella~\ref{tab:winner} riporta la distribuzione dei giudizi
complessivi; la Tabella~\ref{tab:perdomanda} mostra il dettaglio
per ciascuna delle 20 domande.

\begin{table}[tbp]
  \centering
  \caption{Punteggi medi per metrica (scala 1--5, $n = 19$).}
  \label{tab:medie}
  \small
  \renewcommand{\arraystretch}{1.4}
  \begin{tabular}{lcccc}
    \toprule
    \textbf{Metrica} &
    \textbf{Distillation} &
    \textbf{Reflection} &
    \textbf{$\Delta$ (R $-$ D)} \\
    \midrule
    Correctness       & 3.89 & 2.74 & $-$1.16 \\
    Reasoning Quality & 3.74 & 3.05 & $-$0.68 \\
    Clarity           & 4.11 & 2.42 & $-$1.68 \\
    \midrule
    \textbf{Media complessiva} & \textbf{3.91} & \textbf{2.74} & $-$\textbf{1.17} \\
    \bottomrule
  \end{tabular}
\end{table}

\begin{table}[tbp]
  \centering
  \caption{Distribuzione dei giudizi di vittoria su 19 domande valide.}
  \label{tab:winner}
  \small
  \renewcommand{\arraystretch}{1.4}
  \begin{tabular}{lcc}
    \toprule
    \textbf{Esito} & \textbf{Conteggio} & \textbf{Percentuale} \\
    \midrule
    Distillation (A) & 14 & 73.7\% \\
    Reflection (B)   &  5 & 26.3\% \\
    Parità (TIE)     &  0 &  0.0\% \\
    \bottomrule
  \end{tabular}
\end{table}

Il modello di distillation ottiene punteggi superiori in tutte e tre
le dimensioni valutate e vince in 14 domande su 19 (73.7\%), senza
alcun pareggio. Il divario più marcato riguarda la metrica
\emph{Clarity} ($\Delta = -1.68$); il divario minore si osserva sulla
\emph{Reasoning Quality} ($\Delta = -0.68$), suggerendo che il modello
reflection, pur perdendo complessivamente, riesce in alcuni casi a
produrre sequenze di ragionamento di qualità comparabile.

La Tabella~\ref{tab:perdomanda} riporta il dettaglio per domanda.

\begin{table}[tbp]
  \centering
  \caption{Punteggi per domanda. C = Correctness, Q = Reasoning Quality,
    Cl = Clarity. Con * la domanda Q18 per cui il giudice ha restituito
    un errore API.}
  \label{tab:perdomanda}
  \resizebox{\textwidth}{!}{%
  \renewcommand{\arraystretch}{1.2}
  \begin{tabular}{clccccccl}
    \toprule
    \textbf{\#} & \textbf{Domanda (sintesi)} &
    \textbf{D\,C} & \textbf{D\,Q} & \textbf{D\,Cl} &
    \textbf{R\,C} & \textbf{R\,Q} & \textbf{R\,Cl} &
    \textbf{Vincitore} \\
    \midrule
    1  & Resto su acquisto mele        & 5 & 2 & 1 & 5 & 4 & 3 & Reflection \\
    2  & Velocità media treno          & 1 & 1 & 2 & 5 & 5 & 5 & Reflection \\
    3  & Dimensioni rettangolo         & 1 & 1 & 3 & 5 & 5 & 5 & Reflection \\
    4  & Numero ragazzi in classe      & 5 & 5 & 5 & 5 & 3 & 2 & Distillation \\
    5  & Lavoratori e giorni           & 5 & 5 & 5 & 1 & 1 & 1 & Distillation \\
    6  & Prezzo originale camicia      & 5 & 4 & 3 & 1 & 2 & 2 & Distillation \\
    7  & Equazione $x + 3 = 10$        & 5 & 5 & 5 & 4 & 5 & 2 & Distillation \\
    8  & Distanza automobile           & 5 & 5 & 5 & 1 & 2 & 2 & Distillation \\
    9  & Somma angoli pentagono        & 5 & 5 & 5 & 5 & 3 & 2 & Distillation \\
    10 & Divisibile per 4 e 6          & 5 & 4 & 4 & 1 & 2 & 3 & Distillation \\
    11 & Probabilità pallina rossa     & 1 & 2 & 3 & 2 & 4 & 2 & Reflection \\
    12 & Giorno della settimana        & 1 & 1 & 4 & 2 & 4 & 3 & Reflection \\
    13 & Altezza scala appoggiata      & 5 & 5 & 5 & 1 & 2 & 1 & Distillation \\
    14 & Equazione $2(x+4) = 3x-1$     & 5 & 5 & 5 & 1 & 2 & 2 & Distillation \\
    15 & Farina per biscotti           & 5 & 4 & 4 & 1 & 3 & 2 & Distillation \\
    16 & 15\% di 240                   & 1 & 3 & 4 & 1 & 2 & 2 & Distillation \\
    17 & Perimetro quadrato            & 5 & 5 & 5 & 3 & 3 & 2 & Distillation \\
    18 & Vasca con rubinetto*          & {--} & {--} & {--} & {--} & {--} & {--} & Errore API \\
    19 & Sequenza geometrica           & 4 & 4 & 5 & 3 & 2 & 2 & Distillation \\
    20 & Calcolo $f(2)$ con polinomio  & 5 & 5 & 5 & 5 & 4 & 3 & Distillation \\
    \midrule
    \multicolumn{2}{l}{\textbf{Media}} &
    3.89 & 3.74 & 4.11 & 2.74 & 3.05 & 2.42 & \\
    \bottomrule
  \end{tabular}}
\end{table}

\FloatBarrier
% ------------------------------------------------------------
\section{Analisi dei Risultati}
\label{sec:analisi}
% ------------------------------------------------------------

\subsection{Casi in cui la distillation prevale}

Il modello di distillation vince in 14 domande, con un pattern
abbastanza netto: nelle domande in cui la risposta richiede un calcolo
diretto o l'applicazione di una formula elementare (Q4, Q5, Q8, Q9,
Q13, Q14, Q15, Q17), il modello produce risposte corrette, concise
e ben strutturate, che il giudice valuta con il massimo punteggio
su tutte e tre le metriche.

In questi casi il modello reflection produce invece risposte
spesso errate (Correctness $= 1$) o incomplete, con punteggi di
Clarity molto bassi (frequentemente 1 o 2). Questo suggerisce che
il modello da 1B parametri, addestrato su soli 700 campioni, non
ha appreso in modo stabile il formato \texttt{<thinking>}, generando
in molti casi output mal strutturati in cui il ragionamento interno
non converge alla risposta corretta.

Un caso esemplificativo è la domanda Q5 (``Se 5 operai completano
un lavoro in 8 giorni, quanti giorni occorrono a 10 operai?''):
il modello di distillation produce la risposta corretta (4 giorni)
con un ragionamento lineare e punteggi massimi, mentre il modello
reflection ottiene il punteggio minimo su tutte le metriche,
indicando un output probabilmente incompleto o incoerente.

\subsection{Casi in cui la reflection prevale}

Il modello reflection vince in 5 domande (Q1, Q2, Q3, Q11, Q12).
L'analisi di questi casi rivela due sottotipi distinti:

\paragraph{Errori del modello di distillation (Q2, Q3, Q10, Q12).}
In questi casi il modello di distillation produce una risposta errata
(Correctness $= 1$), mentre il modello reflection riesce a fornire
la risposta corretta, probabilmente grazie al ragionamento esplicito
nel blocco \texttt{<thinking>} che funge da scaffolding per il
calcolo. La domanda Q2 (velocità media del treno) è un esempio
rappresentativo: il modello di distillation compie un errore
aritmetico elementare e produce una risposta sbagliata, mentre il
modello reflection percorre correttamente i passaggi di calcolo nel
blocco di thinking e giunge alla risposta esatta.

\paragraph{Parità sulla correttezza, vantaggio sulla qualità del ragionamento (Q1).}
Nella domanda Q1 entrambi i modelli rispondono correttamente, ma
il modello di distillation produce una risposta caratterizzata da
forte ridondanza (lo stesso ragionamento viene ripetuto più volte),
che il giudice penalizza con Clarity $= 1$ e Reasoning Quality $= 2$.
Il modello reflection, pur con Clarity inferiore alla media,
presenta un processo di ragionamento più coerente e viene giudicato
superiore.

\paragraph{Problemi con ragionamento non deterministico (Q11, Q12).}
Per le domande sulla probabilità e sul calendario, entrambi i modelli
mostrano difficoltà (Correctness $\leq 2$ per entrambi), ma il modello
reflection ottiene punteggi di Reasoning Quality superiori (4 vs.\ 1--2),
indicando che anche in presenza di una risposta finale errata il
blocco di thinking mostra un tentativo di ragionamento più articolato.

\subsection{Analisi per tipologia di domanda}

I risultati mostrano una dipendenza dalla tipologia del problema:

\begin{itemize}
  \item \textbf{Aritmetica diretta e algebra elementare} (Q4, Q7, Q8,
    Q9, Q14, Q17, Q20): distillation domina con punteggi massimi.
    Il modello ha appreso efficacemente lo stile conciso di Gemini Flash
    per problemi a soluzione rapida.

  \item \textbf{Problemi con struttura intermedia} (Q2, Q3, Q6, Q13,
    Q15): risultati misti. Distillation vince nella maggioranza, ma
    commette errori su problemi che richiedono più passaggi (Q2, Q3).

  \item \textbf{Ragionamento combinatorio e probabilistico} (Q11, Q12):
    entrambi i modelli faticano, ma reflection mostra un tentativo
    di ragionamento più esplicito, sebbene non sempre corretto.
\end{itemize}

\FloatBarrier
% ------------------------------------------------------------
\section{Discussione}
\label{sec:discussione}
% ------------------------------------------------------------

\subsection{Interpretazione generale}

I risultati indicano che, nelle condizioni sperimentali di questo
lavoro, il modello addestrato con knowledge distillation response-based
da Gemini Flash supera il modello addestrato con reflection tuning
su quasi tutte le metriche e tipologie di domanda. Tuttavia,
prima di trarre conclusioni generali, è necessario contestualizzare
questi risultati rispetto ai vincoli del setup sperimentale.

Il vantaggio della distillation è in larga misura attribuibile alla
\textbf{qualità assoluta delle risposte del teacher}. Gemini Flash
è un modello con capacità significativamente superiori a Llama
3.2 1B, e le sue risposte alle domande di matematica elementare sono
tendenzialmente corrette, concise e ben strutturate. Il modello
student ha potuto imitare questo stile in modo efficace, raggiungendo
punteggi di Clarity molto alti (media 4.11).

Al contrario, il dataset di reflection tuning contiene catene di
pensiero generate da modelli di reasoning di grandi dimensioni,
il cui stile verboso e articolato risulta difficile da replicare
per un modello da 1B parametri con training limitato. Il risultato
è che il modello reflection produce frequentemente output mal
strutturati o incompleti, penalizzati su tutte le metriche.

\subsection{Il ruolo del reasoning esplicito}

Nonostante i risultati complessivamente negativi per il reflection
tuning, i cinque casi in cui questo approccio prevale offrono
un'indicazione interessante: il ragionamento esplicito nel blocco
\texttt{<thinking>} sembra fungere da meccanismo di auto-correzione
su problemi che richiedono più passaggi di calcolo. Quando il modello
di distillation fallisce su questi problemi (Q2, Q3), il modello
reflection riesce a trovare la soluzione corretta percorrendo
esplicitamente i passaggi intermedi.

Questo è coerente con la letteratura sul \emph{chain-of-thought
prompting} \cite{wei2022chain}, che documenta come il
ragionamento esplicito migliori le prestazioni sui problemi che
richiedono composizione di passi. Il limite osservato è che un
modello da 1B parametri non dispone della capacità sufficiente
a mantenere un ragionamento coerente per l'intera durata del blocco
di thinking, degenerando spesso in ripetizioni o divergenze logiche.

\subsection{Confronto con la letteratura}

I risultati sono coerenti con le osservazioni riportate in lavori
recenti sul \emph{small language model fine-tuning}. Lavori come
Alpaca \cite{taori2023alpaca} e Vicuna \cite{chiang2023vicuna}
hanno mostrato che modelli di piccola taglia beneficiano
significativamente dall'imitazione di risposte di teacher più
capaci, anche in assenza di soft targets. D'altra parte, la
superiorità del reflection tuning su modelli di scala maggiore
è documentata in lavori come DeepSeek-R1 \cite{deepseekr1_2025}
e OpenAI o1, che adottano architetture e dataset di training
molto più grandi di quelli disponibili in questo esperimento.

La differenza tra i risultati di questo lavoro e quelli della
letteratura sul reasoning tuning è quindi spiegabile principalmente
dalla scala: con modelli da 7B parametri o superiori e dataset
di training nell'ordine delle decine di migliaia di campioni,
il reflection tuning mostra vantaggi consistenti che a scala 1B
non emergono in modo stabile.

\subsection{Limiti dell'interpretazione}

Alcune cautele sono necessarie nell'interpretare questi risultati:

\begin{itemize}
  \item Il test set è composto da sole 19 domande valide, il che
    non consente di trarre conclusioni statisticamente significative.
    La distribuzione 14--5 osservata potrebbe non essere stabile
    su un insieme di valutazione più ampio.

  \item Il giudice automatico (Gemini Flash) è lo stesso modello
    utilizzato per generare il dataset di distillation. Non si può
    escludere un bias stilistico che favorisce risposte simili a
    quelle che il giudice produce autonomamente.

  \item Le 20 domande coprono esclusivamente matematica e logica
    elementare. I risultati potrebbero differire su domini diversi,
    in particolare su problemi che richiedono ragionamento multi-passo
    più complesso, dove il blocco \texttt{<thinking>} potrebbe
    offrire un vantaggio maggiore.
\end{itemize}

